\documentclass[aspectratio=169]{beamer}

\makeatletter
\def\input@path{{template/}}
\makeatother

% Switch between modes:
% \usetheme[dark]{uob}	(default)	dark title + dark slides
% \usetheme[light]{uob}				light title + light slides
% \usetheme[panda]{uob}				dark title + light slides
% \usetheme[zebra]{uob}				light title + dark slides
\usetheme[dark]{uob}

\usepackage{booktabs}

\title{University of Birmingham Beamer Template}
\date{\today}
\author{Joseph Chamberlain}

\begin{document}

\begin{frame}
    \titlepage
\end{frame}

\section{You can declare sections like this.}

\begin{frame}{Blocks make things stand out.}
    \begin{block}{Block}
        This is the default block. It uses the University's preferred yellow colours. 
    \end{block}
    
    \begin{alertblock}{Alert Block}
        This is an alert block. It uses the University's red colours.
    \end{alertblock}
    
    \begin{examples}
        This is for examples. It uses the University's blue colours.
    \end{examples}
\end{frame}

\begin{frame}{You can put maths in.}
    \begin{equation}
        a^2 = \alpha^2 + c^2
    \end{equation}
\end{frame}

\section{Another section here}

\begin{frame}{Figures look rather nice.}
    \begin{figure}
        \centering
        % Comment / uncomment the appropriate logo for the color theme
        \includegraphics[width=0.5\linewidth]{template/logos/UoB_dark.png}
        %\includegraphics[width=0.5\linewidth]{template/logos/UoB_light.png}
        \caption{Caption}
        \label{fig:enter-label}
    \end{figure}
\end{frame}

\end{document}
